\chapter{Fetal Alcohol Spectrum Disorders (FASD) Screening}
\section*{What Is This Test?} FASD screening identifies children or adults who may need assessment for effects associated with prenatal alcohol exposure; there is no single diagnostic laboratory test.
\section*{Alternative Names} FASD assessment; fetal alcohol exposure screening; prenatal alcohol exposure evaluation.
\section*{Principle} Developmental, behavioral, learning, growth, facial, neurologic, and medical information is combined with exposure history when available.
\section*{Why This Test Is Ordered} It is considered for developmental delay, learning or attention problems, behavior concerns, characteristic growth or facial findings, or known prenatal alcohol exposure.
\section*{Primary vs Secondary Test} Screening is preliminary; multidisciplinary diagnostic evaluation is required and should not depend on a perfect exposure history.
\section*{How This Test Is Performed} Clinicians review records and examine development, cognition, behavior, growth, and physical findings across settings.
\section*{What Preparation Is Needed} Bring prenatal, birth, school, therapy, and medical records when available.
\section*{Understanding Results} A screen indicates need for evaluation, not a diagnosis. Findings may overlap with other developmental conditions.
\section*{Follow-up Tests} Neuropsychological, speech-language, hearing, vision, genetic, neurologic, and behavioral assessments may follow.
\section*{References} \begin{enumerate}\item MedlinePlus. Fetal Alcohol Spectrum Disorders Screening.\item Centers for Disease Control and Prevention. FASD diagnosis and prevention resources.\end{enumerate}
\clearpage\section*{Understanding Results and Follow-up}
The laboratory report should identify the specimen, method, units, reference interval or decision limit, and whether the result is positive, negative, abnormal, or inconclusive. Reference intervals vary with age, sex, pregnancy, specimen type, assay, and laboratory; a result outside a range is not automatically a diagnosis, and a result within a range does not always exclude disease. Discuss unexpected results with the ordering clinician, who can decide whether repeat or confirmatory testing, treatment, or referral is appropriate. Do not change medicines or delay urgent care based on an isolated result.
\section*{Regulatory Status and Authorization Date}This chapter describes a test category, not a specific commercial product. FDA, CDSCO, EU IVDR, and MHRA approval or registration dates therefore cannot be assigned without the assay manufacturer, product name, instrument, and intended use. For laboratory-developed tests, an individual agency approval date may not exist. See the Preface for the interpretation of regulatory status.
\clearpage
